\section{Hecke modifications and transverse cotangent sections}
\label{sec:hecke}

We now turn the Lie-theoretic data of Proposition~\ref{lie:input} into a
modification of a bundle and its Higgs field.  The moving point of the
modification produces a nonzero covector on the curve.  Two further
properties will be needed: at fixed input, and at fixed output, the
corresponding relative cotangent section has an invertible derivative.
In the next section, the fixed-output derivative will isolate the
distinguished stalk. The fixed-input derivative will complete the
coordinates for its quadratic local model.

\subsection{Integral correspondences and the open Satake kernel}

Keep the trait $S$, curve $\mathscr X/S$, and bundle stack $\mathcal B$
from Section~\ref{sec:specialization}.  A one-point Hecke modification is a
quadruple $(P,P',x,\alpha)$, where $x$ is a point of $\mathscr X$ and
$\alpha$ identifies $P'$ with $P$ away from its graph.  Write
\[
 p(P,P',x,\alpha)=P,\qquad o(P,P',x,\alpha)=P',\qquad
 \ell_H(P,P',x,\alpha)=x.
\]
We fix the following relative-position convention.  In a formal coordinate
$t$ at $x$ and an input frame, modification by a cocharacter $\lambda$
means that the output frame is given by $\lambda(t)$.  Its relative
position depends only on the Weyl orbit of $\lambda$.  Let
$\mu=\lambda^+$ be the dominant representative and put
\[
 m=\langle 2\rho,\mu\rangle,
 \qquad \mu^\vee=-w_0\mu,
 \qquad \mathcal Z=\mathcal B\times_S\mathscr X,
\]
where $2\rho$ is the sum of the positive roots and $w_0$ is the longest
Weyl element.  Denote the exact-position correspondence by $H=H^\lambda$,
and set
\[
 \widetilde p=(p,\ell_H):H\longrightarrow\mathcal Z,
 \qquad q=(o,\ell_H):H\longrightarrow\mathcal Z.
\]

\begin{lemma}[Hecke models and kernels]\label{hecke:models}
Both $\widetilde p$ and $q$ are representable, separated, and smooth of
relative dimension $m$.  The correspondence $H$ is open in a bounded
correspondence $\overline H$ for which both projections to $\mathcal Z$
are representable and proper.  On the output side the bound has dominant
label $\mu^\vee$.

For the irreducible dual-group representation $V$ corresponding to this
bound, with the input and output convention just specified, the generic
Hecke functor has the form
\begin{equation}\label{hecke:kernel}
 \mathrm H_V(F)=\overline q_{K,*}\mathcal I_K.
\end{equation}
On $H_K$, the integrand $\mathcal I_K$ is $p_K^*F$ tensored with a fixed
invertible constant complex.  On a fixed-leg affine-Grassmannian fiber,
the perverse intersection-cohomology normalization of that complex is
$E[m]$, with the constant Tate line if weight normalization is used.
\end{lemma}

\begin{proof}
Use the split integral affine Grassmannian, whose loops are
$G(A((t)))$ and whose positive loops are $G(A[[t]])$.  Its integral
Schubert bound is projective, and the positive-loop orbit is smooth,
open, and fiberwise dense, of relative dimension $m$; see
\cite[\S\S1.2 and~2.1]{RicharzScholbach}.  We use these models by base change from
$\mathbb Z$.  The variable $t$ is the curve parameter, including when the
base trait has mixed characteristic.

A coordinate along the graph of $x$ and a frame on the formal disc
identify modifications of a fixed input bundle with the corresponding
Grassmannian spaces.  These descriptions glue by
Beauville--Laszlo descent \cite{BeauvilleLaszlo}.  More explicitly, the
graph is an effective Cartier divisor.  The gluing theorem for modules
also glues the coordinate algebra, multiplication, and coaction of an
affine $G$-torsor; the torsor identity can be checked on the two gluing
pieces.  A disc frame exists locally: lift a frame on the graph
successively through its infinitesimal neighborhoods using smoothness of
$G$.  Changes of frame and formal coordinate preserve the orbit and its
bound.  Consequently the properness, smoothness, and dimensions descend
to the input-and-leg projection.

Swapping the bundles and inverting $\alpha$ is an involution of the Hecke
correspondence.  On double cosets it replaces $\mu$ by $-\mu$, whose
dominant representative is $\mu^\vee$.  It carries the bound to the
oppositely labelled bound.  This proves the assertions for $q$, since
$\langle 2\rho,\mu^\vee\rangle=m$.  Separatedness of the open
correspondence follows from that of the proper bound.  This involution
is used on the correspondence, rather than as a purported inversion map
on the right-quotient affine Grassmannian.

The fibers here include an identification of the fixed bundle with the
specified bundle.  They are relative-position spaces: an automorphism
of the moving bundle compatible with its identification off the graph
is the identity.  In particular, no further quotient by automorphisms
of the fixed bundle is taken.

For the last assertion, geometric Satake gives the intersection complex
of the bound.  Its restriction to the smooth open orbit is the shifted
constant sheaf, with the indicated optional Tate normalization; see
\cite[Proposition~4.1]{Richarz}.  The Hecke construction tensors this kernel with the ordinary pullback
$p_K^*F$ and pushes forward along the proper output-and-leg projection;
see \cite[\S\S4.2.1--4.2.4]{GR}.  Relabelling by $\mu^\vee$ if necessary
matches these projections with the chosen representation.  A fixed
overall shift or Tate line in the Hecke normalization changes neither
this description nor any vanishing assertion below.  Only the open-orbit
restriction is being identified; the kernel on the boundary remains its
intersection complex.
\end{proof}

\subsection{The transverse modification}

On the special fiber, a cotangent vector at $P\in\mathcal B_s(k)$ means
an element of $H^0$ of the cotangent-complex fiber. Deformation theory
and vector-bundle Serre duality \cite[Tag~0BS4]{Stacks}, followed by
the invariant form $\kappa$, identify this
space with
\[
 T_P^*\mathcal B_s
  =H^1(X,\mathfrak g_P)^*
  =H^0(X,\mathfrak g_P\otimes\omega_X),
\]
where $\mathfrak g_P$ is the adjoint bundle.  We call such a vector $A$
a Higgs field.  For a representable smooth map, pullback on these
cotangent spaces is injective and its image is the kernel of passage to
the relative cotangent.  This follows from the cotangent triangle; it
does not assert that the map to relative cotangents is surjective for
maps between stacks.

\begin{proposition}[Transverse modification]\label{hecke:transverse}
Assume the hypotheses of Proposition~\ref{lie:input}.  Let
$P\in\mathcal B_s(k)$ and let $A\in T_P^*\mathcal B_s$ be non-nilpotent.
There exist a point $x\in X(k)$, a cocharacter $\lambda$, a modification
$h\in H^\lambda_s(k)$ from $P$ to a bundle $P'$, and a covector
\[
 \theta=(A',\xi)\in T^*_{(P',x)}(\mathcal B_s\times X),
 \qquad \xi\ne0,
\]
such that $A'$ agrees with $A$ away from $x$ and
\begin{equation}\label{hecke:cotangent-identity}
 dp_h^*A=dq_h^*\theta.
\end{equation}
Moreover, form the fixed-side fibers
\[
 H^{\mathrm{in}}_{P,x}=\widetilde p^{-1}(P,x),
 \qquad H^{\mathrm{out}}_{P',x}=q^{-1}(P',x),
\]
with the fixed-bundle identifications retained.  On the first fiber,
restriction of $p^*A$ to the $q$-relative tangent bundle defines a section
$S_A$ of $\Omega_{H_s/\mathcal Z_s,q}$.  On the second fiber, restriction
of $o^*A'$ to the $\widetilde p$-relative tangent bundle similarly defines
$S_{A'}$.  Both sections vanish at $h$, and both derivatives
\begin{align*}
 (dS_A)_h &:T_hH^{\mathrm{in}}_{P,x}
                  \longrightarrow (T_hH^{\mathrm{out}}_{P',x})^*,\\
 (dS_{A'})_h &:T_hH^{\mathrm{out}}_{P',x}
                  \longrightarrow (T_hH^{\mathrm{in}}_{P,x})^*
\end{align*}
are isomorphisms.  In particular, each section has an isolated reduced
zero at $h$.
\end{proposition}

\begin{proof}
Since $A$ is non-nilpotent, choose $x\in X(k)$ where its value is
non-nilpotent. Choose a formal coordinate $t$ and a bundle frame at $x$,
and write $A=a(t)\,dt$. Chevalley restriction identifies the nilpotent
fiber with $\chi^{-1}(\chi(0))$, so $\chi(a(0))\ne\chi(0)$.
Proposition~\ref{lie:input} then supplies a formal change of frame,
a Levi subgroup $M$, and a cocharacter
$\lambda$ of $Z(M)$ such that
\begin{equation}\label{hecke:lie-data}
 a(t)\in\mathfrak m[[t]],\qquad
 \operatorname{ad}(a(0))
       \text{ is invertible on }\mathfrak g/\mathfrak m,
 \qquad \kappa(a(0),d\lambda)\ne0.
\end{equation}
Here $d\lambda$ is the value of the cocharacter differential on the
canonical generator of $\Lie(\mathbb G_m)$.  After constant conjugation,
we regard $\lambda$ as an element of the fixed cocharacter lattice; it
therefore also labels the integral model of Lemma~\ref{hecke:models}.
Modify by $\lambda(t)$.  Since $\lambda$ centralizes $M$, the transported
field is regular in the output frame and still has expression
$a(t)\,dt$.  Together with $A$ away from $x$, it defines the global Higgs
field $A'$ on $P'$.  This central-in-a-Levi construction and the resulting
leg covector are the modification used in
\cite[\S\S20.4--20.6]{AGKRRV}.  We prove the derivative assertions
explicitly.

\smallskip\noindent
\emph{The two tangent spaces.}
In the input frame put
\[
 L=\mathfrak g[[t]],\qquad
 L'=\Ad(\lambda(t))L,\qquad I=L\cap L'.
\]
Let $\mathfrak g_j$ be the weight-$j$ space for the algebraic
$\mathbb G_m$-action defined by $\lambda$.  The integers $j$ denote
characters of $\mathbb G_m$, without reduction modulo $p$.  Then
\begin{equation}\label{hecke:weight-quotients}
 \begin{split}
 L'&=\bigoplus_j t^j\mathfrak g_j[[t]],\\
 P_1:=L/I&=\bigoplus_{j>0}
                   \mathfrak g_j\otimes_k k[[t]]/(t^j),\\
 P_2:=L'/I&=\bigoplus_{j<0}
                   \mathfrak g_j\otimes_k(t^jk[[t]]/k[[t]]).
 \end{split}
\end{equation}
An infinitesimal disc automorphism $1+\epsilon D$, $D\in L$, acts on the
output lattice while fixing $(P,x)$.  Its stabilizer tangent is $I$.
The orbit is smooth, and the root formula gives
$\dim_k(L/I)=m$, so this computes its full tangent.  The same argument
with the output fixed gives
\begin{equation}\label{hecke:tangents}
 T_hH^{\mathrm{in}}_{P,x}=P_1,
 \qquad T_hH^{\mathrm{out}}_{P',x}=P_2.
\end{equation}
The notation $1+\epsilon D$ uses the canonical identification of the
kernel over dual numbers with the Lie algebra; no exponential map is
required.

There is a perfect pairing
\begin{equation}\label{hecke:residue-pairing}
 \beta:P_1\times P_2\longrightarrow k,
 \qquad \beta(D,C)=\Res_{t=0}\kappa(D,C)\,dt.
\end{equation}
Indeed, $L$ is its own annihilator for the residue pairing on
$\mathfrak g((t))$, and invariance of $\kappa$ gives the same assertion
for $L'$.  Thus the pairing descends to the displayed quotients.
More explicitly, $\kappa$ pairs $\mathfrak g_j$ perfectly with
$\mathfrak g_{-j}$; the coefficient of $t^a$ in the first space, for
$0\le a<j$, pairs with the coefficient of $t^{-a-1}$ in the second.
This proves perfection, including when some $j$ is divisible by $p$.

\smallskip\noindent
\emph{Transport of the covector and motion of the leg.}
Choose the gluing sign convention in which an infinitesimal transition
$1+\epsilon C$ gives the boundary class of the principal part $C$.
Serre duality pairs this class with $A$ by
$\Res\kappa(a,C)\,dt$.  For an output-fixed tangent $C\in L'$, this
residue is zero, since $a\in L'$ and $L'$ is self-annihilating.
Consequently $dp_h^*A$ vanishes in the $q$-relative cotangent, so it is
$dq_h^*\theta$ for a unique covector $\theta$.

Its bundle component is $A'$.  To check this, choose a point $x'\ne x$
and vary both bundle transitions there while leaving the modification
at $x$ unchanged.  Principal parts at $x'$ surject onto
$H^1(X,\mathfrak g_{P'})$: for sufficiently large $N$, Serre vanishing
gives $H^1(X,\mathfrak g_{P'}(Nx'))=0$.  The pairings on these
directions agree because $A=A'$ near $x'$.  Write therefore
$\theta=(A',\xi)$.

Now keep the input bundle fixed and move the leg from $t=0$ to
$t=\epsilon$.  The output transition becomes $g_\epsilon=\lambda(t-\epsilon)$,
and its first-order displacement is
\[
 (\delta g)g^{-1}=-\frac{d\lambda}{t}.
\]
Evaluating the cotangent identity on this motion yields
\[
 0=-\kappa(a(0),d\lambda)+\xi(\partial_t),
 \qquad
 \xi=\kappa(a(0),d\lambda)\,dt\big|_x\ne0.
\]
This is also the residue formula of
\cite[Equation~(20.20)]{AGKRRV}.  Reversing the gluing boundary convention
would reverse the common sign and would leave all subsequent assertions
unchanged.

\smallskip\noindent
\emph{The first relative derivative.}
The section $S_A$ is defined along the entire fixed-input fiber, using
its specified identification with $P$.  Its value at $h$ is zero by the
preceding calculation.  For $D\in L$, deform the output lattice to
\[
 L'_\epsilon=\Ad(1+\epsilon D)L'.
\]
A test vector $C\in L'$ extends as
$C_\epsilon=\Ad(1+\epsilon D)C$ in the deformed lattice.  The input
field remains the fixed $a(t)\,dt$.  Differentiating the residue
pairing gives
\begin{equation}\label{hecke:derivative}
 \begin{split}
 \big\langle(dS_A)_h(D\bmod I),C\bmod I\big\rangle
  &=\Res\kappa(a,[D,C])\,dt\\
  &=\Res\kappa([a,D],C)\,dt
   =\beta(\operatorname{ad}(a)D,C).
 \end{split}
\end{equation}
Because $S_A(h)=0$, this derivative is independent of how the test
vector is extended.  Also $\operatorname{ad}(a)$ preserves $L$, $L'$,
and $I$, so the formula is independent of representatives.

It remains to prove invertibility.  The operator
$\operatorname{ad}(a(t))$ preserves every weight summand in
\eqref{hecke:weight-quotients}.  Since $\lambda$ is central in $M$,
each nonzero weight space injects as an invariant summand into
$\mathfrak g/\mathfrak m$.  Its endomorphism induced by
$\operatorname{ad}(a(0))$ is therefore invertible by
\eqref{hecke:lie-data}.  The determinant of the corresponding
$k[[t]]$-linear block is a unit.  It follows that
$\operatorname{ad}(a(t))$ is invertible on every truncated summand of
$P_1$ and $P_2$.  Equation~\eqref{hecke:derivative}, together with the
perfect pairing~\eqref{hecke:residue-pairing}, now proves that
$(dS_A)_h:P_1\to P_2^*$ is an isomorphism.  No regularity of $\lambda$
outside $M$ is required; any additional zero-weight directions do not
occur in these quotients.

\smallskip\noindent
\emph{The reverse derivative.}
The value of $S_{A'}$ at $h$ is zero because $a\in L$ and $L$ is
self-annihilating.  Keeping the output fixed, vary the input lattice by
$C\in L'$ and extend a test vector $D\in L$ to
$\Ad(1+\epsilon C)D$.  The same calculation, with the same gluing
convention, gives
\[
 \big\langle(dS_{A'})_h(C\bmod I),D\bmod I\big\rangle
   =\Res\kappa(a,[C,D])\,dt
   =\Res\kappa([a,C],D)\,dt.
\]
Invertibility on $P_2$ and residue duality prove the second assertion.
Both sections are sections of rank-$m$ vector bundles on smooth
$m$-dimensional fibers.  Their invertible derivatives make their zeros
at $h$ reduced and isolated.

When $m=0$, both tangent spaces are zero and both derivatives are
isomorphisms between zero spaces.  The smooth zero-dimensional fibers
already isolate $h$, and the moving-leg computation still applies.
Thus the conclusion includes central modifications and torus factors.
\end{proof}

The proposition gives both kinds of control needed below: the nonzero
leg covector makes a suitable hypersurface smooth over the output
bundle, while the two relative derivatives isolate one point and produce
the quadratic coordinates around it.
